\subsection{g-模的特征标}

设 \(\Delta\) 是用加法记号书写的交换幺半群，\(\mathbf{Z}[\Delta] = \mathbf{Z}^{(\Delta)}\) 是幺半群 \(\Delta\) 在 \(\mathbf{Z}\) 上的代数（《代数》，第 III 章，§2，第 6 节）。以 \((e^{\lambda})_{\lambda \in \Delta}\) 表示 \(\mathbf{Z}[\Delta]\) 的典则基。对于所有 \(\lambda, \mu \in \Delta\)，我们有 \(e^{\lambda + \mu} = e^{\lambda}e^{\mu}\)。若 0 是 \(\Delta\) 的中性元，则 \(e^0\) 是 \(\mathbf{Z}[\Delta]\) 的单位元；记为 1。

设 E 是 \(\Delta\)-分次的、域 \(\kappa\) 上的向量空间，并令 \((\mathrm{E}^{\lambda})_{\lambda\in\Delta}\) 为其分次。若每个 \(E^{\lambda}\) 都是有限维的，则 E 的特征标（记为 \(\operatorname{ch}(\operatorname{E})\)）是由 \((\dim\operatorname{E}^{\lambda})_{\lambda\in\Delta}\) 给出的 \(Z^{\Delta}\) 中的元素。若 E 本身是有限维的，

\[
\operatorname{ch} (\mathrm{E}) = \sum_ {\lambda \in \Delta} (\dim \mathrm{E} ^ {\lambda}) e ^ {\lambda} \in \mathbf {Z} [ \Delta ].\tag{5}
\]

设 \(E', E, E''\) 是 \(\Delta\)-分次向量空间，使得 \(E'^{\lambda}, E^{\lambda}, E''^{\lambda}\) 关于 \(\kappa\) 都是有限维的，并且 \(0 \longrightarrow E' \longrightarrow E \longrightarrow E'' \longrightarrow 0\) 是次数为 0 的分次同态的正合序列。立即可得

\[
\operatorname{ch} (\mathrm{E}) = \operatorname{ch} (\mathrm{E} ^ {\prime}) + \operatorname{ch} (\mathrm{E} ^ {\prime \prime}).\tag{6}
\]

特别地，若 \(F_{1}, F_{2}\) 是 \(\Delta\)-分次向量空间，且 \(F_{1}^{\lambda}\) 与 \(F_{2}^{\lambda}\) 关于 \(\kappa\) 都是有限维的，则

\[
\operatorname{ch} \left(\mathrm{F} _ {1} \oplus \mathrm{F} _ {2}\right) = \operatorname{ch} \left(\mathrm{F} _ {1}\right) + \operatorname{ch} \left(\mathrm{F} _ {2}\right).\tag{7}
\]

若 \(F_{1}\) 和 \(F_{2}\) 是有限维的，则还有

\[
\operatorname{ch} \left(\mathrm{F} _ {1} \otimes \mathrm{F} _ {2}\right) = \operatorname{ch} \left(\mathrm{F} _ {1}\right). \operatorname{ch} \left(\mathrm{F} _ {2}\right).\tag{8}
\]

例。假设 \(\Delta = N\)。设 T 是一个不定元。存在唯一的同构，从代数 Z{[}N{]} 到代数 Z{[}T{]}，将 \(e^{n}\) 映到 \(T^{n}\)（对所有 \(n \in N\)）。对于任意有限维 N-分次向量空间 E，\(\operatorname{ch}(E)\) 在 Z{[}T{]} 中的像是 E 的 Poincaré 多项式（第 V 章，§5，第 1 节）。

设 E 是满足 \(E = \sum_{\lambda \in h^{*}} E^{\lambda}\) 且每个 \(E^{\lambda}\) 都是有限维的 g-模。我们知道 \((\mathrm{E}^{\lambda})_{\lambda \in \mathfrak{h}^{*}}\) 是向量空间 E 的一个分次。下文中，我们保留记号 \(\operatorname{ch}(\mathrm{E})\)，用来表示把 E 看作 \(h^{*}\)-分次向量空间时的特征标。因此，特征标 \(\operatorname{ch}(\mathrm{E})\) 是 \(Z^{h^{*}}\) 中的元素。若 E 是有限维的，则 \(\operatorname{ch}(\mathrm{E}) \in \mathbf{Z}[\mathrm{P}]\)。根据公式 (6) 和第 6 节的引理 3，存在唯一的同态，从群 \(\mathcal{R}(\mathfrak{g})\) 到 Z{[}P{]}，对任意有限维 g-模 E 将其映到 \(\operatorname{ch}(\mathrm{E})\)；这个同态记为 ch。关系 (8) 表明 ch 是从环 \(\mathcal{R}(\mathfrak{g})\) 到环 Z{[}P{]} 的环同态。

备注。P 的每个元素定义一个维数为 1 的单 h-模，因而给出从群 Z{[}P{]} 到群 \(\mathcal{R}(\mathfrak{h})\) 的同态；该同态是环的单射同态。显然复合

\[
\mathcal {R} (\mathfrak {g}) \longrightarrow \mathbf {Z} [ \mathrm{P} ] \longrightarrow \mathcal {R} (\mathfrak {h})
\]

是由 \(\mathfrak{h}\) 到 \(\mathfrak{g}\) 的典则嵌入定义的同态（第 6 节）。

Weyl 群 W 通过自同构作用于群 P，因而作用于 \(Z^{P}\)。对于所有 \(\lambda \in P\) 和所有 \(w \in W\)，我们有 \(we^{\lambda} = e^{w\lambda}\)。设 \(Z[P]^{W}\) 是 Z{[}P{]} 的子环，由在 W 下不变的元素组成。

引理 6. 若 \(\lambda \in \mathrm{P}_{++}\)，则 \(\operatorname{ch}[\lambda] \in \mathbf{Z}[\mathrm{P}]^{\mathrm{W}}\)。\(\operatorname{ch}[\lambda]\) 的唯一极大项（第 VI 章，§3，第 2 节，定义 1）是 \(e^{\lambda}\)。

第一个断言由第 1 节命题 2 的推论 2 得出，第二个断言由 §6，第 1 节，命题 1 (ii) 得出。

定理 2. (i) 设 \((\varpi_{\alpha})_{\alpha \in \mathrm{B}}\) 是关于 B 的基本权族。设 \((\mathrm{T}_{\alpha})_{\alpha \in \mathrm{B}}\) 是不定元族。由 \(f\mapsto f(\left([ \varpi_{\alpha} ]\right)_{\alpha \in \mathrm{B}})\) 定义的从 \(\mathbf{Z}[(\mathrm{T}_{\alpha})_{\alpha \in \mathrm{B}}]\) 到 \(\mathcal{R}(\mathfrak{g})\) 的映射是环同构。

\begin{enumerate}
\def\labelenumi{(\roman{enumi})}
\setcounter{enumi}{1}
\item
  由 \(\operatorname{ch}\) 诱导的从 \(\mathcal{R}(\mathfrak{g})\) 到 \(\mathbf{Z}[\mathrm{P}]\) 的同态，诱导出从环 \(\mathcal{R}(\mathfrak{g})\) 到环 \(\mathbf{Z}[\mathrm{P}]^{\mathrm{W}}\) 的同构。
\item
  设 E 是有限维 g-模。若 \(\operatorname{ch} E = \sum_{\lambda \in P_{++}} m_{\lambda} \operatorname{ch}[\lambda]\)，则 E 中最高权为 \(\lambda\) 的同构型分支的长度为 \(m_{\lambda}\)。
\end{enumerate}

族 \(([\lambda])_{\lambda \in \mathrm{P}_{++}}\) 是 \(\mathbf{Z}\)-模 \(\mathcal{R}(\mathfrak{g})\) 的一个基，而族 \((\operatorname{ch}[\lambda])_{\lambda \in \mathrm{P}_{++}}\) 是 \(\mathbf{Z}\)-模 \(\mathbf{Z}[\mathrm{P}]^{\mathrm{W}}\) 的一个基（引理 6，以及第 VI 章，§3，第 4 节，命题 3）。这证明了 (ii) 和 (iii)。断言 (i) 由 (ii)、引理 6 以及第 VI 章，§3，第 4 节，定理 1 得出。

推论。设 E、E' 是有限维 g-模。则 E 同构于 E'，当且仅当 ch E = ch E'。

这由定理 2 (ii) 以及 E、E' 是半单模这一事实得出。

